% page 182 du fac-simile (image p-181 du scan)
% bloc 162.6 x 229.0 mm ; marge gauche 8.0 mm
\pgc{-12.700mm}{0.000mm}{
\l{\sou{fonteno}. Aparato qua varsas, per robineto, la aquo ye qua}
\l{\cel{3}onu plenigis lu o quan onu adduktas per tubi. - DEFIR.}
\l{}
\l{\sou{for}.\cel{2}Prepoziciono qua signifikas eskarteso, plu o min gram-}
\l{\cel{3}da, de... - EF.}
\l{}
\l{\sou{forco}. (\sou{fiziko)}. I. La kauzo di movo. - II. Kauzo nekonocata}
\l{\cel{3}di ula movi di materio (gravito, kaloro, elektro, e c.).}
\l{\cel{3}EFIS.}
\l{}
\l{\sou{forcato.}\cel{2}Persono qua subisas puniso afliktiva ed infamiganta}
\l{\cel{3}qua konsistis, olim, ye remar sur la galeri rejala ; nun,}
\l{\cel{3}ye deportesar en remparo-cirkuito fortifikita. - FIS.}
\l{}
\l{\sou{forcepso}. (\sou{medic.)}\cel{2}Instrumento ek fero, du-brancha quan uzas}
\l{\cel{3}la akushisto en la kazi di desfacileso, por sizar du-latere}
\l{\cel{3}la kapo di la infanti ed atraktar lu ad su. - DEFIS.}
\l{}
\l{\sou{forejo}.\cel{2}Feno di la prati naturala od artificala, herbi, veje-}
\l{\cel{3}tanti qui uzesas por nutrar la brutaro. - DEFIRS.}
\l{}
\l{\sou{foresto}. Vasta surfaco, areo, de tereno arboroza. - DEFIS.}
\l{}
\l{\sou{forfeto}. (\sou{yuro-cienco}). Kontrato inter du personi, ek qui ica}
\l{\cel{3}su obligas prenor - ita, livror, po preco determinita anti-}
\l{\cel{3}cipe, ula laboruri, ula servi. - FI.}
\l{}
\l{\sou{forfikulo}. (\sou{zool.}) Insekto ortoptera, di qua la ventro finas}
\l{\cel{3}per du hoketi pincatra. - L.\cel{1}\sou{forficula}. - IL.}
\l{}
\l{\sou{forjar}.\cel{2}(\sou{trans.)}\cel{2}Laborar (metalo) sur amboso, per fairo e}
\l{\cel{3}martelo. - EFS.}
\l{}
\l{\sou{forko}. I. Instrumento kun longa mancho ligna, finanta per du}
\l{\cel{3}o tri branchi pintoza, ek ligno o fero (por inversige}
\l{\cel{3}agitar feno, palio, movar sterko, deprenar o kargar garbi).}
\l{\cel{3}- II. Objekto qua kelke similesas forko. - DEFI.}
\l{}
\l{\sou{formo}.\cel{2}I.Aspekto extera di korpo. - II. (\sou{metaf}.) Karaktero}
\l{\cel{3}sub qua kozo su manifestas. - III. Maniero extera agar segun}
\l{\cel{3}ula reguli.- IV. Maniero expresar penso. - DEFIRS.}
\l{}
\l{\sou{formacar}. (\sou{trans.}) Facar (ulu, ulo) donante a lu la eso e la}
\l{\cel{3}formo. - DEFIRS.}
\l{}
\l{\sou{formalino.}\cel{1}(\sou{kemio}).H. CO2H. - DEF.}
\l{}
\l{\sou{formato}. Dimensiono quan onu donas a libro, a jurnalo, segun}
\l{\cel{3}ke onu faldas la imprimo-folio unfoye, dufoye, trifoye.DFIRS.}
\l{}
\l{\sou{formiko}. (\sou{zool.)}\cel{2}Mikra insekto ek la klaso \textquotedbl{}himenopteri\textquotedbl{}, qua}
\l{\cel{3}facas por su (en la kavaji di la olda arbori) lojeyo ube lu}
\l{\cel{3}vivas societane. - L.\cel{1}\sou{formica}.\cel{2}FIS.}
}
